\begingroup
\section{Node DF0A: intrinsic variance and the far stratum}
\label{module:DF0A}
\noindent\textit{Audited source: \nolinkurl{DF0_variance.tex}.}
\subsection{The intrinsic equality-stratum variance}

Let
\begin{equation}\label{DF0A:eq:data}
 h_i>0,\qquad \sum_i h_i=2\pi=:L,
 \qquad a=\frac{2\pi}{N},\qquad
 S=\max_i h_i+a.
\end{equation}
The quadratic cell bubble and its sampled Hilbert conjugate are
\begin{align}
 q_\alpha(\theta_i+x)
 &=\frac12\min\{x^2,(h_i-x)^2\},\qquad 0\le x\le h_i,
 \label{DF0A:eq:q}\\
 T(\alpha)&=\sum_i h_i|\Hil q_\alpha(c_i)|^2,
 \qquad c_i=\theta_i+h_i/2.\label{DF0A:eq:T}
\end{align}
Put $A_k=\sum_i h_i^k$ and define
\begin{equation}\label{DF0A:eq:D}
 \ell^2:=\frac{A_3}{L},
 \qquad
 \boxed{\mathfrak D(\alpha)
 :=\sum_i h_i(h_i^2-\ell^2)^2
 =A_5-\frac{A_3^2}{L}.}
\end{equation}

\begin{lemma}[Pair-variance formula and equality strata]
\label{DF0A:lem:pair}
One has
\begin{equation}\label{DF0A:eq:pair}
 \mathfrak D(\alpha)
 =\frac1{2L}\sum_{i,j}h_ih_j(h_i^2-h_j^2)^2.
\end{equation}
Consequently $\mathfrak D\ge0$, and it vanishes exactly when all
positive cell lengths are equal.  After zero-cell deletion, these are
precisely the collapsed-uniform strata.
\end{lemma}

\begin{proof}
Expand the double sum and use $\sum_i h_i=L$ and
$\sum_i h_i^3=A_3$.  Equality in a weighted variance holds precisely
when $h_i^2$ is constant on its positive support.
\end{proof}

Recall the quartic Bregman charge
\begin{equation}\label{DF0A:eq:J}
 \Jfour(\alpha)=\sum_i(h_i-a)^2(h_i^2+2ah_i+3a^2).
\end{equation}

\begin{proposition}[The intrinsic variance is paid by the DAG charge]
\label{DF0A:prop:D-J}
For every positive fan,
\begin{equation}\label{DF0A:eq:D-J}
 \boxed{\mathfrak D(\alpha)\le S\Jfour(\alpha).}
\end{equation}
The estimate is independent of the minimum cell and all adjacent ratios.
\end{proposition}

\begin{proof}
The number $\ell^2=A_3/L$ is the weighted least-squares minimizer of
$m\mapsto\sum_i h_i(h_i^2-m)^2$.  Hence, using the competitor $m=a^2$,
\begin{align*}
 \mathfrak D(\alpha)
 &\le\sum_i h_i(h_i^2-a^2)^2\\
 &=\sum_i h_i(h_i-a)^2(h_i+a)^2\\
 &\le S\sum_i(h_i-a)^2(h_i^2+2ah_i+3a^2)
 =S\Jfour(\alpha).
\end{align*}
\end{proof}

\subsection{Exact relation with the cubic null form}

Stage 144 proves
\begin{align}
 \Cnull(q_\alpha^3)&=V(\alpha)-T(\alpha),\label{DF0A:eq:C-VT}\\
 V(\alpha)&=\frac{A_5}{320}-\frac{A_3^2}{1152\pi}.
 \label{DF0A:eq:V}
\end{align}

\begin{proposition}[Intrinsic defect identity]
\label{DF0A:prop:defect}
For every positive fan,
\begin{equation}\label{DF0A:eq:defect}
 \boxed{
 \Cnull(q_\alpha^3)-\frac1{720}\sum_i h_i^5
 =\frac1{576}\mathfrak D(\alpha)-T(\alpha).}
\end{equation}
In particular the sharp inequality
\begin{equation}\label{DF0A:eq:sharp-candidate}
 \Cnull(q_\alpha^3)\ge\frac1{720}\sum_i h_i^5
\end{equation}
is equivalent to $T\le\mathfrak D/576$.
\end{proposition}

\begin{proof}
Since $L=2\pi$,
\[
 \frac1{320}-\frac1{720}=\frac1{576},
 \qquad
 \frac{A_3^2}{1152\pi}=\frac1{576}\frac{A_3^2}{L}.
\]
Insert these two equalities in \eqref{DF0A:eq:C-VT}--\eqref{DF0A:eq:V}.
\end{proof}

\begin{remark}[Logical status of the sharp constant]
Equation \eqref{DF0A:eq:defect} is proved; inequality
\eqref{DF0A:eq:sharp-candidate} is not used as an established fact.  The
unconditional proof only needs an unspecified universal constant in the
following smaller endpoint leaf.
\end{remark}

\begin{obligation}[AVS: adaptive equality-stratum variance sampling]
\label{DF0A:obl:AVS}
Prove
\begin{equation}\label{DF0A:eq:AVS}
 \boxed{T(\alpha)\le C\mathfrak D(\alpha)}
\end{equation}
for every positive fine fan.  By Proposition~\ref{DF0A:prop:D-J}, AVS implies
HQ and hence DF.  Unlike the deleted HMS/QPS route, AVS vanishes on every
collapsed-uniform endpoint stratum and contains no path velocity.
\end{obligation}

\subsection{A global absolute sampling bound}

We use the following mesh-free sampling theorem in the precise special
case needed here.

\begin{lemma}[Adaptive truncated-Hilbert sampling]
\label{DF0A:lem:sampling}
Let $\{I_i\}$ be any interval partition of the circle, with lengths
$h_i$ and midpoints $c_i$.  If $f\in L^2(\T)$ satisfies
\begin{equation}\label{DF0A:eq:self-zero}
 \operatorname {pv}\int_{I_i}\frac1{2\pi}
 \cot\frac{c_i-y}{2}\,f(y)\,dy=0
 \qquad\text{for every }i,
\end{equation}
then
\begin{equation}\label{DF0A:eq:sampling}
 \sum_i h_i|\Hil f(c_i)|^2\le C\|f\|_2^2.
\end{equation}
The constant is independent of all interval ratios.
\end{lemma}

\begin{proof}
The left side of \eqref{DF0A:eq:self-zero} is the part removed by symmetric
truncation at radius $h_i/2$.  For every $x\in I_i$, the
scale-restricted Cotlar inequality gives
\[
 |\Hil f(c_i)|\le C\{M(\Hil f)(x)+Mf(x)\}.
\]
Indeed, split at radius $h_i$; moving the centre from $c_i$ to $x$
changes the far part by a Hardy--Littlewood average, and the intervening
annulus is another such average.  Square, integrate on $I_i$, sum the
disjoint cells, and use the $L^2$ boundedness of $M$ and $\Hil$.
\end{proof}

\begin{theorem}[Unconditional absolute endpoint bound]
\label{DF0A:thm:absolute}
For every positive fan,
\begin{equation}\label{DF0A:eq:absolute}
 \boxed{T(\alpha)\le C A_5=C\sum_i h_i^5.}
\end{equation}
\end{theorem}

\begin{proof}
On its own cell, $q_\alpha$ is even about $c_i$, while the Hilbert kernel
is odd.  Thus \eqref{DF0A:eq:self-zero} holds with $f=q_\alpha$.  Moreover
direct integration of the two half-bubbles gives
\[
 \|q_\alpha\|_2^2=\sum_i\frac{h_i^5}{320}.
\]
Apply Lemma~\ref{DF0A:lem:sampling}.
\end{proof}

\subsection{All fans away from the equality strata are closed}

\begin{theorem}[Far-stratum AVS and HQ]
\label{DF0A:thm:far}
Fix $0<\delta<1$.  If
\begin{equation}\label{DF0A:eq:far}
 \mathfrak D(\alpha)\ge\delta A_5,
\end{equation}
then
\begin{equation}\label{DF0A:eq:far-conclusion}
 T(\alpha)\le C_\delta\mathfrak D(\alpha)
 \le C_\delta S\Jfour(\alpha).
\end{equation}
Thus AVS, HQ, and the pure cubic row DF are unconditional on the complete
far-stratum region, including arbitrary minimum cells and ratios.
\end{theorem}

\begin{proof}
Theorem~\ref{DF0A:thm:absolute} and \eqref{DF0A:eq:far} give the first inequality;
Proposition~\ref{DF0A:prop:D-J} gives the second.  Stage 144 then gives
HQ $\Rightarrow$ DF.
\end{proof}

\subsection{Quantitative selection of the nearest collapsed-uniform stratum}

It remains to treat $\mathfrak D<\delta A_5$.  This condition has a
strong geometric consequence which we now make exact.

\begin{lemma}[Concentration of physical angular mass]
\label{DF0A:lem:concentration}
Assume $0<\delta<1$ and
$\mathfrak D<\delta A_5$.  Then
\begin{equation}\label{DF0A:eq:D-small}
 \mathfrak D\le\frac{\delta}{1-\delta}L\ell^4.
\end{equation}
For $0<\rho<1$ let
\begin{equation}\label{DF0A:eq:bad}
 \mathcal B_\rho=\{i:|h_i/\ell-1|\ge\rho\},
 \qquad B_\rho=\sum_{i\in\mathcal B_\rho}h_i.
\end{equation}
Then
\begin{equation}\label{DF0A:eq:bad-mass}
 \boxed{
 \frac{B_\rho}{L}
 \le\frac{\delta}{(1-\delta)\rho^2}.}
\end{equation}
\end{lemma}

\begin{proof}
Equation \eqref{DF0A:eq:D} gives $A_5=\mathfrak D+L\ell^4$;
this proves \eqref{DF0A:eq:D-small}.  On a bad cell,
$|h_i^2-\ell^2|\ge\rho\ell^2$.  Hence
$\rho^2\ell^4B_\rho\le\mathfrak D$, and
\eqref{DF0A:eq:bad-mass} follows.
\end{proof}

Let $\mathcal G_\rho$ be the complementary good set and
$M=|\mathcal G_\rho|$.  If
$\delta<(1-\delta)\rho^2$, Lemma~\ref{DF0A:lem:concentration} ensures
$M\ge1$.  Define the collapsed-uniform reference fan $\beta$ by
\begin{equation}\label{DF0A:eq:beta}
 \beta_i=\begin{cases}b:=L/M,&i\in\mathcal G_\rho,\\0,&i\in\mathcal B_\rho.
 \end{cases}
\end{equation}
After deleting its zero cells, $\beta$ is the regular $M$-fan; therefore
all its Hilbert midpoint samples vanish.

\begin{lemma}[Quantitative distance to the selected stratum]
\label{DF0A:lem:reference}
Put $\varepsilon=B_\rho/L$.  Then
\begin{equation}\label{DF0A:eq:b-ell}
 \frac{1-\rho}{1-\varepsilon}
 \le\frac b\ell\le
 \frac{1+\rho}{1-\varepsilon}.
\end{equation}
If $\rho+\varepsilon\le1/4$, then
\begin{equation}\label{DF0A:eq:reference-charge}
 \boxed{
 b^3\sum_{i\in\mathcal G_\rho}(h_i-b)^2
 +\sum_{i\in\mathcal B_\rho}h_i(h_i^2-\ell^2)^2
 \le C_\rho\mathfrak D(\alpha).}
\end{equation}
\end{lemma}

\begin{proof}
Since every good length lies between $(1-\rho)\ell$ and
$(1+\rho)\ell$,
\[
 M(1-\rho)\ell\le L-B_\rho
 \le M(1+\rho)\ell.
\]
Using $b=L/M$ gives \eqref{DF0A:eq:b-ell}.

Write $e_i=h_i-\ell$ on the good set.  There
\[
 h_i(h_i^2-\ell^2)^2\ge c_\rho\ell^3e_i^2.
\]
Also
\[
 M(b-\ell)=\sum_{i\in\mathcal G_\rho}e_i+B_\rho.
\]
Therefore Cauchy--Schwarz, $M\asymp L/\ell$, and
$\rho^2\ell^4B_\rho\le\mathfrak D$ give
\begin{align*}
 b^3\sum_{\mathcal G_\rho}(h_i-b)^2
 &\le C_\rho\ell^3\sum_{\mathcal G_\rho}e_i^2
   +C_\rho\frac{\ell^3B_\rho^2}{M}\\
 &\le C_\rho\mathfrak D
   +C_\rho\ell^4\frac{B_\rho^2}{L}
 \le C_\rho\mathfrak D.
\end{align*}
The bad term in \eqref{DF0A:eq:reference-charge} is a sub-sum of
$\mathfrak D$.
\end{proof}

The preceding lemmas reduce the near-stratum region to the following
strictly localized endpoint estimate.

\begin{obligation}[ELCA: equality-stratum first-separation square]
\label{DF0A:obl:ELCA}
For the good/bad decomposition and collapsed-uniform reference
\eqref{DF0A:eq:beta}, prove directly from the endpoint kernel-difference
identity of stage 152 that
\begin{multline}\label{DF0A:eq:ELCA}
 T(\alpha)\\
 \le C_\rho\left\{
 b^3\sum_{i\in\mathcal G_\rho}(h_i-b)^2
 +\sum_{i\in\mathcal B_\rho}h_i(h_i^2-\ell^2)^2
 \right\}.
\end{multline}
The regular $M$-fan is used only after zero-cell deletion.  A maximal
consecutive bad block is kept intact until its first separating scale;
the two endpoint kernels may not be estimated separately inside that
block.  Good--good rectangles use the signed variable-grid discrete
Calder\'on--Zygmund square, while every rectangle meeting a bad block is
paid by the second term in \eqref{DF0A:eq:ELCA}.  Descendants are squared
before ancestors are summed.
\end{obligation}

\begin{proposition}[ELCA is the only remaining pure endpoint row]
\label{DF0A:prop:ELCA}
ELCA implies AVS in the near-stratum region.  Together with
Theorem~\ref{DF0A:thm:far}, it implies AVS, HQ, and DF for every positive fan.
\end{proposition}

\begin{proof}
Lemma~\ref{DF0A:lem:reference} bounds the right side of \eqref{DF0A:eq:ELCA} by
$C_\rho\mathfrak D$.  This is AVS.  The far and near regions exhaust all
fans.  Proposition~\ref{DF0A:prop:D-J} and stage 144 then give HQ and DF.
\end{proof}
\endgroup
